\chapter*{Declarations}
\addcontentsline{toc}{chapter}{Declarations}

\section*{Use of Generative AI}
I acknowledge the use of Claude Code and the \texttt{claude\_agent\_sdk} (Anthropic,
\url{https://claude.com}) throughout the development of the evaluation subsystem
described in this report, including for code scaffolding, debugging and iterative
refactoring. This is consistent with the subject matter of the project itself, which builds evaluation tooling for Claude
Skills. I also acknowledge the use of Claude (Anthropic, \url{https://claude.ai}) to generate an outline for background study. . I confirm that the analysis, conclusions and technical content presented in this report are my own work.

% TODO (please confirm before submission): exact tool versions used, any other tools
% (e.g. ChatGPT, GitHub Copilot) not listed here, and whether the Department requires
% a more detailed log of prompts used for the report-writing assistance above.

\section*{Ethical Considerations}
This project builds internal developer tooling, an evaluation subsystem for Claude
Skills, within Man Group's existing Skill Manager application. It does not involve
human participants, personal data, or any data collected specifically for the
purposes of this research. The test prompts and quality rubrics used to exercise
the system during development were either synthetic examples or generated by the
project's own AI-assisted test-suggestion feature, not real user or client data.
All data touched by the system, namely skill definitions and evaluation run
results, is Man Group's existing internal engineering data, held under the firm's
existing information-security and access-control policies, and no new categories of sensitive data were introduced by this project.

Ethical considerations were discussed with my supervisor at the outset of the
project. Given the absence of human subjects, personal data or any external
dataset, no formal College ethical approval was required or sought. I am not
aware of any ethical issues arising from this project.

% TODO (please confirm before submission): that this matches the outcome of the
% ethics review conducted with your supervisor; add an approval reference number
% here if one was issued.

\section*{Sustainability}
Because each skill execution under test spawns a full Node.js and Claude Code
subprocess, compute and memory efficiency were treated as a first-class
engineering constraint rather than an afterthought, with a direct environmental
benefit. Skill execution defaults to the smaller, cheaper Claude Haiku 4.5 model
for the frequent iteration loop an author works in, reserving the more capable,
and more compute-intensive, Claude Sonnet 4.6 for the comparatively infrequent
judging step. Judge calls are routed through Man Group's lightweight
\texttt{man.ai} HTTP gateway rather than the heavyweight SDK execution path,
avoiding an unnecessary subprocess and its associated resource footprint. A
pod-wide concurrency semaphore caps the number of simultaneous subprocesses no
matter how many evaluation runs fan out, and per-prompt and whole-run timeouts
bound any runaway compute. The system's memory footprint was measured
empirically, rather than provisioned by guesswork, so the production deployment
uses no more cluster capacity than the measured requirement plus a modest safety
margin. Together, these choices minimise the compute, and therefore the energy,
consumed both during development and in ongoing production use.

\section*{Availability of Data and Materials}
The source code for this project is held in Man Group's internal Forgejo git
service, as part of the existing Skill Manager repository. This repository is
private to Man Group and cannot be made publicly available due to commercial
confidentiality. No separate public repository, dataset or standalone web
application is provided; all information necessary to understand the system's
design, implementation and evaluation is contained within this report. Access to
the private repository can be arranged for examination purposes on request,
subject to Man Group's approval.

% TODO (please confirm before submission): whether examiners will actually be
% offered repo access, and update this paragraph (or point to a shared drive or
% demo environment instead) if not.
